\section{Sleep Loss, Stimulants, and Catastrophic Executive Breakdown}\label{sec:executive_breakdown}
Sleep deprivation alone induces striking dissociations: vigilance and motor speed can stay intact while goal switching,
selective attention, and metacognitive monitoring collapse \citep{couyoumdjian2010, slama2018, garcia2021}. Stimulants intensify the
split. D-amphetamine restores subjective alertness and arithmetic accuracy during the first hours of deprivation, yet verbal
reasoning and planning lag far behind \citep{newhouse1989}. Higher doses or stacked compounds (tyrosine, phentermine, caffeine) help
with working memory span but do little for executive reconfiguration \citep{magill2003}. Animal work reveals a mechanistic reason: the
combination of sleep loss and amphetamine accelerates sensitization and drug seeking, meaning dopamine pathways are being driven
harder even as receptors downregulate \citep{berro2018, saito2014}. Human PET imaging confirms the receptor story directly: total sleep
deprivation lowers striatal D2/D3 availability, and the magnitude of that drop predicts failures in attention and working memory
\citep{volkow2012, seeman2011}. In OBH terms, stimulants dial up the precision of whatever model is already running, so a deprived
brain becomes an overconfident but brittle predictor. Executive collapse is therefore not a mystery of willpower but a measurable
consequence of violating the offline cycle the system was built to depend on.
